
\begin{figure}[t]
    {\small
    \textbf{Pattern normal form} (\rocqlink{https://edgardschi.github.io/mailbox-types-rocq/MailboxTypes.MailboxPatterns.html\#PNFLit}) \hfill \fbox{$E \pnf F$}
    \begin{center}
          \begin{prooftree}
            \infer0{E \pnf \mbzero}
          \end{prooftree}
          \qquad
          \begin{prooftree}
            \infer0{E \pnf \mbone}
          \end{prooftree}
          \qquad
          \begin{prooftree}
            \hypo{\mbpres{E}{m}{E'}}
            \hypo{F \mbpeq E'}
            \infer2{E \pnf \mbmsg{m} \mbcomp F}
          \end{prooftree}
          \qquad
          \begin{prooftree}
            \hypo{E \pnf F_1}
            \hypo{E \pnf F_2}
            \infer2{E \pnf F_1 \mbchoice F_2}
          \end{prooftree}
      \end{center}
  }
    \vspace{-1em}

\caption{Pattern normal form.}
\label{fig:pnf}
\end{figure}

\subsection{Pattern Normal Form}
When reasoning about the typing of guards, both the Mailbox Calculus and Pat
require the type of a mailbox to be in \emph{pattern normal form} (PNF), which
is a particular structure of pattern that reflects the subpatterns expected by
each guard clause.
%
Specifically, a pattern $E$ is in PNF if it is a sum of subpatterns of the form
$\mbzero$ (for a
$\keyword{fail}$ guard), $\mbone$ (for a $\keyword{free}$ guard) and $\mbmsg{m}
\mbcomp F$ (for a $\keyword{receive}$ guard), where $F$ is equivalent to the
residual of $E$ with respect to $\msg{m}$. This allows a correct type to be
given to a re-bound mailbox name after receiving.

Figure~\ref{fig:pnf} shows the formal definitions for PNF; we write $\pnf E$ for
$E \pnf E$.
%
The following example derivation shows that
pattern $(\mbmsg{m} \mbcomp \mbone) \mbchoice \mbone$ is in PNF:

\vspace{-0.3cm}
\begin{center}
\scalebox{0.80}{
\begin{prooftree}
  \infer0{\mbpres{\mbmsg{m}}{m}{\mbone}}
  \infer0{\mbpres{\mbone}{m}{\mbzero}}
  \infer2{\mbpres{(\mbmsg{m} \mbcomp \mbone) \mbchoice \mbone}{m}{(\mbone \mbcomp \mbone) \mbchoice (\mbmsg{m} \mbcomp \mbzero)}}
  \hypo{\mbone \mbpeq (\mbone \mbcomp \mbone) \mbchoice (\mbmsg{m} \mbcomp \mbzero)}
  \infer2{(\mbmsg{m} \mbcomp \mbone) \mbchoice \mbone  \pnf (\mbmsg{m} \mbcomp \mbone)}
  \infer0{(\mbmsg{m} \mbcomp \mbone) \mbchoice \mbone  \pnf \mbone}
  \infer2{(\mbmsg{m} \mbcomp \mbone)\mbchoice \mbone \pnf (\mbmsg{m} \mbcomp \mbone) \mbchoice \mbone}
\end{prooftree}
}
\end{center}
\vspace{-0.7cm}



\subsection{Typing Rules}
\label{sec:typing-rules}
We now have all the necessary definitions for specifying the typing rules of Pat.
These rules are similar to the original ones by Fowler et al.
\cite{fowlerSpecialDeliveryProgramming2023}, with a few modifications to fit our
definitions of variables, environments and terms.

\begin{figure}[!t]
    \footnotesize
  \textbf{Typing rules for programs and functions}
  (\rocqlink{https://edgardschi.github.io/mailbox-types-rocq/MailboxTypes.TypingRules.html\#WellTypedProgram})
  (\rocqlink{https://edgardschi.github.io/mailbox-types-rocq/MailboxTypes.TypingRules.html\#WellTypedDefinition}) \;
  \hfill \fbox{$\WellTypedProg$} \fbox{$\WellTypedDef{\FunctionDef}$}
\begin{mathpar}
\inferrule
{
\mathcal{P} = (\mathcal{S}, \mathcal{D}, t_0) \\
\forall \defname .\ \WellTypedDef{\mathcal{D}(\defname)} \\
\WellTypedTerm{\emptyset}{t_0}{\unittype}
}
{\WellTypedProg}

\inferrule*
{
\WellTypedTerm{A, \emptyset}{t}{B}
}
{\WellTypedDef{\FunctionDef}}

\end{mathpar}

\textbf{Typing rules for terms} (\rocqlink{https://edgardschi.github.io/mailbox-types-rocq/MailboxTypes.TypingRules.html\#WellTypedTerm})
\hfill
\fbox{$\WellTypedTermP{\Gamma}{\mathcal{P}}{t}{A}$}
%
\begin{mathpar}
\inferrule[\VAR]
{
\EmptyEnv{\Gamma}
}
{\WellTypedTerm{\envInsert{x}{A}{\Gamma}}{\Var{x}}{A}}

\inferrule[Bool]
{
\EmptyEnv{\Gamma}
}
{\WellTypedTerm{\Gamma}{b}{\booltype}}

\inferrule[\UNIT]
{
\EmptyEnv{\Gamma}
}
{\WellTypedTerm{\Gamma}{()}{\unittype}}

\inferrule[\APP]
{
\mathcal{D}_{\mathcal{P}}(\defname)
 = \keyword{def} \ \defname(A) : B \ \{t\} \\\\
\WellTypedTerm{\Gamma}{v}{A}
}
{\WellTypedTerm{\Gamma}{\defname(v)}{B}}

\inferrule[\LET]
{
\envcomb{\Gamma_1}{\Gamma_2}{\Gamma} \\
\WellTypedTerm{\Gamma_1}{t_1}{\returnableop{A}} \\\\
\WellTypedTerm{\envInsert{0}{\returnableop{A}}{\Gamma_2}}{t_2}{B}
}
{\WellTypedTerm{\Gamma}{\TLet{t_1}{t_2}}{B}}
%
\quad
%
\inferrule[\NEW]
{
\EmptyEnv{\Gamma}
}
{\WellTypedTerm{\Gamma}{\TNew}{\mbtin{\mbone}^\ureturnable}}
%
\quad
%
\inferrule[\SEND]
{
\mathcal{S}_{\mathcal{P}}(\msg{m}) = T \\
\WellTypedTerm{\Gamma_1}{v_1}{\mbtout{\mbmsg{m}^\usecondclass}} \\\\
\WellTypedTerm{\Gamma_2}{v_2}{\secondclassop{T}} \\
\disenvcomb{\Gamma_1}{\Gamma_2}{\Gamma}
}
{\WellTypedTerm{\Gamma}{\TSend{v_1}{m}{v_2}}{\unittype}}

\inferrule[\GUARD]
{
E \mbpinc F \\
\pnf F \\
\WellTypedTerm{\Gamma_1}{v}{\mbtin{F^\ureturnable}} \\\\
\WellTypedGuards{\Gamma_2}{\overrightarrow{g}}{A}{F} \\
\disenvcomb{\Gamma_1}{\Gamma_2}{\Gamma}
}
{\WellTypedTerm{\Gamma}{\TGuard{v}{E}{\overrightarrow{g}}}{A}}

\inferrule[\SUB]
{
\Gamma \subtype \Gamma' \\\\
A \subtype B \\
\WellTypedTerm{\Gamma'}{t}{A}
}
{\WellTypedTerm{\Gamma}{t}{B}}

\inferrule[\SPAWN]
{
\secondclassop{\Gamma} = \Gamma' \\\\
\WellTypedTerm{\Gamma'}{t}{\unittype}
}
{\WellTypedTerm{\Gamma}{\TSpawn{t}}{\unittype}}
\end{mathpar}

  \textbf{Typing rules for guard sequences} (\rocqlink{https://edgardschi.github.io/mailbox-types-rocq/MailboxTypes.TypingRules.html\#WellTypedGuards})
  \hfill \fbox{$\WellTypedGuardP{\Gamma}{\overrightarrow{g}}{A}{E}$}
  \begin{center}
    \begin{prooftree}
      \hypo{\WellTypedGuard{\Gamma}{g}{A}{E}}
      \infer1[$\SINGLE$]{\WellTypedGuards{\Gamma}{\{g\}}{A}{E}}
    \end{prooftree}
    \qquad
    \begin{prooftree}
      \hypo{\WellTypedGuard{\Gamma}{g}{A}{E}}
      \hypo{\WellTypedGuards{\Gamma}{\overrightarrow{g_s}}{A}{E'}}
      \infer2[$\SEQ$]{\WellTypedGuards{\Gamma}{\{g, \overrightarrow{g_s}\}}{A}{E \mbchoice E'}}
    \end{prooftree}
  \end{center}
  \textbf{Typing rules for guards} (\rocqlink{https://edgardschi.github.io/mailbox-types-rocq/MailboxTypes.TypingRules.html\#WellTypedGuard})
  \hfill \fbox{$\WellTypedGuardP{\Gamma}{g}{A}{E}$}
    \begin{mathpar}
        \inferrule
        [Fail]
        { }
        {\WellTypedGuard{\Gamma}{\GFail}{A}{\mbzero}}

        \inferrule
        [Free]
        { \WellTypedTerm{\Gamma}{t}{A} }
        { \WellTypedGuard{\Gamma}{\GFree{t}}{A}{\mbone} }

        \inferrule
        [Receive]
        { \mathcal{S}_{\mathcal{P}}(\msg{m}) = T \\
          \mathsf{base(T)} \lor \mathsf{base(\Gamma)} \\\\
          \WellTypedTerm
            {\envInsert{0}{\secondclassop{T}}{(\envInsert{0}{\mbtin{E}^\ureturnable}{\Gamma})}}
            {t}
            {B}
        }
        { \WellTypedGuard{\Gamma}{\GReceive{m}{t}}{B}{\mbmsg{m} \mbcomp E}}
    \end{mathpar}

  \caption{Pat typing rules}
  \label{fig:typing-rules}
\end{figure}

Figure \ref{fig:typing-rules} shows the typing rules for Pat, parameterised by a
program $\mathcal{P}$; we omit $\mathcal{P}$ in the rules for readability.
We highlight some of the typing rules.
Whenever an empty environment is required, we use predicate $\EmptyEnv{\Gamma}$ instead of checking for the empty environment $\emptyset$.
The insertion operation $\envInsert{x}{A}{\Gamma}$ ensures that $\Gamma$ contains type $A$ at position $x$.

The $\SEND$ rule types a send expression $\TSend{v_1}{m}{v_2}$, where message $\msg{m}$ with payload $v_2$ is sent to mailbox $v_1$.
The mailbox associated with $v_1$ has to be of type $\mbtout{\mbmsg{m}^\usecondclass}$, i.e.\ it must be capable of sending the message $\msg{m}$.
Payload $v_2$ is required to be of the type given by the message's signature with second-class usage.
Both derivations for $v_1$ and $v_2$ must occur under separate environments $\Gamma_1$ and $\Gamma_2$, i.e.\ they are only share base types and their combination yields $\Gamma$.

% The typing rule for a guard expression $\TGuard{v}{E}{\overrightarrow{g}}$ ensures that the pattern $E$ is a subpattern of some pattern $F$ in PNF.
% Additionally, mailbox $v$ must have type $\mbtin{F^\ureturnable}$ under environment $\Gamma_1$, ensuring that it is capable of receiving messages according to $F$.
% The requirement for the returnable usage avoids use-after-free errors.

Rule $\RECEIVE$ types the receive guard $\GReceive{m}{t}$, retrieving message $\msg{m}$ from the mailbox, with type $B$ and pattern $\mbmsg{m} \mbcomp E$.
Either $\Gamma$ must only consist of base types or the payload's type $T$ is a base type.
This is a syntactic restriction, ensuring that no existing mailbox name is shadowed by a message's payload.
The continuation $t$ must be typeable under $\Gamma$ extended by $\mbtin{E^\ureturnable}$ and $\secondclassop{T}$.
Type $\mbtin{E^\ureturnable}$ represents the mailbox after $\msg{m}$ is consumed.
Because of the usage of de Bruijn indices, the order in which the types are inserted into the environment is important.
Since inserting a type at position $0$ corresponds to simply adding it to the front of the environment,
the environment $\envInsert{0}{\secondclassop{T}}{(\envInsert{0}{\mbtin{E}^\ureturnable}{\Gamma})}$ is equal to $\secondclassop{T}, \mbtin{E}^\ureturnable, \Gamma$.
Thus, a variable $0$ would map to type $\secondclassop{T}$ and a variable $1$ would map to type $\mbtin{E}^\ureturnable$.

The typing relation for a sequence of guards enables us to
typecheck each guard individually, while also ensuring the correct pattern
structure of the overall expression.  
In the original definition of typing for guard sequences, Fowler et al. \cite{fowlerSpecialDeliveryProgramming2023} provide only a single rule of the following form:

{\small
    \vspace{-0.5em}
\begin{center}
  \begin{prooftree}
    \hypo{(\WellTypedGuard{\Gamma}{g_i}{A}{E_i})_i}
    \infer1{\WellTypedGuards{\Gamma}{\overrightarrow{g}}{A}{E_1 \mbchoice \dots \mbchoice E_n}}
  \end{prooftree}
\end{center}
}

By splitting this single rule into two distinct ones ($\SINGLE$ and $\SEQ$), we
provide a recursive definition for typing of sequences.  This avoids using lists
of well-typed relations as a premise in the rule, and in turn, Rocq is able to
automatically come up with appropriate inductive schemes.

\ifappendix
An example of a type derivation in Pat can be found in Appendix \ref{sec:type-checking-example}.
\fi

% SJF: I think this is way more detail than necessary.
%   Rule $\SINGLE$ types a sequence $\{g\}$, consisting of a single guard, by
%   requiring $g$ to be typable using the typing rules for guards.  The other rule,
%   $\SEQ$, types a sequence $\{g, \overrightarrow{g_s}\}$ consisting of a guard $g$
%   followed by a sequence of guards $\overrightarrow{g_s}$.  If guard $g$ is
%   typable with type $A$ and pattern $E$, then $\overrightarrow{g_s}$ must be
%   typable under the same environment with the same type and some pattern $E'$.
%   The pattern of $\{g, \overrightarrow{g_s}\}$ is the combination of both patterns
%   by choice: $E \mbchoice E'$.  This way, for a sequence of guards $\{g_1, g_2,
%   \dots, g_n\}$, its pattern is required to be of the form $E_1 \mbchoice E_2
%   \mbchoice \dots \mbchoice E_n$, corresponding to the notion that any available
%   guard within a guard expression can be triggered.
%

\paragraph{Mechanization.}
The mechanization of Pat's typing rules closely follows the pen-and-paper definitions.
However, proving properties about well-typed terms requires technical lemmas, such as inversion lemmas or
canonical forms of terms.
For example, we had to prove that every well-typed Boolean value under environment $\Gamma$ implies that
$\Gamma$ is a subenvironment of an empty environment (\rocqlink{https://edgardschi.github.io/mailbox-types-rocq/MailboxTypes.TypingRules.html\#weak_BTBool_2}).
These types of proofs would usually be omitted in classical pen-and-paper proofs, but in our case,
even these kinds of proofs required technical lemmas such as properties about environment subtyping.
